\section{Il Portale incontra Parco Caprotti}\label{sec:caprotti}

Parco Caprotti è il luogo nominato nella fonte originale.
Ripartire da qui permette di mettere il progetto in relazione
con uno spazio preciso, con la sua storia e con le persone
che lo attraversano. La collocazione conta quanto l'oggetto:
cambia il modo di avvicinarsi, il tempo che si può dedicare
alla musica, la possibilità di fermarsi accanto a qualcuno.

Visit Bergamo riconduce palazzo e parco alla famiglia Caprotti,
industriali brianzoli, e alla trasformazione ottocentesca del
luogo \cite{visit_protestanti}. Questa storia offre il primo
contesto della ricerca. Va approfondita attraverso documenti
e archivio, per comprendere come il giardino sia cambiato
e quali usi abbiano accompagnato quei cambiamenti.

\subsection{Il Padiglione come luogo possibile}
Il Padiglione del Tè porta nella proposta un rapporto
particolare tra interno ed esterno. La presenza del vetro
suggerisce un'attività che può essere vista prima di entrarvi:
qualcuno si ferma nel giardino, osserva una persona alla
tastiera e riconosce un invito possibile.

Immagino Il Portale dentro questo rapporto, con una
collocazione che lasci leggibile l'architettura e conservi
la relazione con il parco. La forma contemporanea
dell'oggetto e quella del Padiglione dovrebbero trovare
una misura comune. Per capire quale, occorre disegnarle,
confrontarle e provarle insieme.

La visibilità richiede sensibilità. Per qualcuno sarà un
invito, per altri un motivo d'esitazione. Lo spazio dovrebbe
permettere anche di ascoltare prima di partecipare, o di
provare senza sentirsi al centro di una piccola platea.
Questo è uno dei punti in cui architettura ed educazione
si toccano concretamente.

\subsection{Un precedente da cui imparare}
Il resoconto di Prima Bergamo del febbraio 2024 descrive
l'attività di Contemporary Locus nel parco e nel Padiglione
durante il 2023. Riferisce anche di un questionario preparato
con l'Orto Botanico e di un'analisi degli esiti condotta
con l'Università degli Studi di Bergamo
\cite{primabergamo2024}.

Il precedente suggerisce un modo di procedere: ascoltare
chi frequenta il luogo e mettere in relazione ricerca
culturale e usi quotidiani. Per Il Portale sarebbe prezioso
capire chi si ferma, chi passa oltre, chi vorrebbe suonare
e chi teme che l'attività cambi il carattere del giardino.
Gli enti citati appartengono a quel precedente; il loro
eventuale coinvolgimento nel nostro progetto va ancora costruito.

\subsection{Una presenza musicale nella vita del parco}
L'attrattività che cerco è legata a ciò che accade:
una persona prova, due persone si ascoltano, un piccolo
gruppo costruisce una sequenza. L'oggetto può destare
curiosità con la sua forma; la qualità dell'esperienza
può dare una ragione per fermarsi e ritornare.

Si potrebbero immaginare momenti di esplorazione libera,
incontri presidiati da musicisti e brevi attività
intergenerazionali. Una raccolta di materiali del
territorio potrebbe offrire percorsi verso scuole,
associazioni e altri luoghi musicali. È una proposta
di programmazione da discutere con possibili interlocutori,
a partire dalle risorse effettivamente disponibili.

Il suono dovrà convivere con il riposo, le conversazioni
e il silenzio di chi frequenta il parco. Orari, modalità
d'ascolto e presidio fanno parte della composizione
di questa esperienza pubblica. La qualità civica del
progetto dipenderà anche da quanto saprà rispettare
le persone che scelgono di non usarlo.
